\clearpage
\fancyhead[L]{\small BEOBACHTUNGEN.}
\SourceLeaf{468}
\begin{center}
\rule{0.55\linewidth}{0.4pt}\\[2pt]
Astronomische Nachrichten. Band X. Nr. 225. Seite 143. 1832 Juli 3.\\[2pt]
\rule{0.55\linewidth}{0.4pt}
\end{center}

\begin{center}[\textsc{Gauss} an \textsc{Schumacher.}]\quad Göttingen 1832 Mai 12.\end{center}

{\rightskip=0pt plus 1em Bei dem letzten Mercursdurchgang habe ich am Eintritt die erste Berührung durch eine Wolke, die zweite durch einen Zufall verloren. Den Austritt hingegen habe ich beobachtet\par}

\begin{center}
\begin{tabular}{ll}
innere Berührung bei dem Austritt \mbox{.\,.\,.\,.\,.\,.\,.} & $4^{\rm u}25^{\rm m}32^{\rm s}1$ \\
äussere Berührung \mbox{.\,.\,.\,.\,.\,.\,.\,.\,.\,.\,.\,.\,.\,.} & $4\;28\;22.5$ \\
\end{tabular}
\end{center}

Ausserdem habe ich seine Culmination am Meridiankreise beobachtet

\medskip
\noindent Mai 5. AR.\quad $2^{\rm u}49^{\rm m}50^{\rm s}36$\par
\noindent $\delta+16^\circ28'43^{\rm s}6$ der Nullpunkt aus dem Nadir bestimmt.\par
\noindent\hspace*{1em}---\quad ---\quad $44.4$ der Nullpunkt aus $\alpha$ Can. min. mit $\alpha$ Decl. aus \textsc{Enckes} \mbox{I. B.}

\clearpage
\hypertarget{native-468-pallas}{}
\begin{center}
\rule{0.55\linewidth}{0.4pt}\\[2pt]
Astronomische Nachrichten. Band XI. Nr. 262. Seite 403. 1834 Juli 9.\\[2pt]
\rule{0.55\linewidth}{0.4pt}
\end{center}

\begin{center}Opposition der \emph{Pallas} 1834.\end{center}

Von Herrn Hofrath \textsc{Gauss} in Göttingen habe ich folgende von ihm am Meridiankreise gemachte Beobachtungen der Pallas erhalten. Wir haben hier keine erhalten.

\begin{center}
\begin{tabular}{l|l|l}
1834 Febr. 6 & $8^{\rm u}41^{\rm m}47^{\rm s}02$ & $-20^\circ44'23''7$ \\
\hfill 9 & $8\;39\;39.10$ & $-19\;41\;52.9$ \\
\hfill 10 & $8\;38\;57.80$ & $-19\;20\;0.3$ \\
\end{tabular}
\end{center}
\begin{flushright}\textsc{Schumacher.}\end{flushright}

\clearpage
\hypertarget{native-468-occultation}{}
\begin{center}
\rule{0.55\linewidth}{0.4pt}\\[2pt]
Astronomische Nachrichten. Band XVI. Nr. 361. Seite 5. 1838 October 18.\\[2pt]
\rule{0.55\linewidth}{0.4pt}
\end{center}

\begin{center}\emph{Sternbedeckung} beobachtet auf der Göttinger Sternwarte.\end{center}

\begin{center}
\begin{tabular}{llll}
Eintritt $\chi$ Leonis 1838 Juni 27. & $10^{\rm u}9^{\rm m}17^{\rm s}7$ & M. Z. & \textsc{Gauss.} \\
 & $10\;9\;16.9$ & --- & \textsc{Goldschmidt.} \\
\end{tabular}
\end{center}

Durch die Güte des Herrn Hofraths und Ritters \textsc{Gauss} mitgetheilt.
\begin{flushright}\textsc{Schumacher.}\end{flushright}
